\documentclass[12pt]{article}

\usepackage{amsmath, amssymb, amsthm}
\usepackage{geometry}
\usepackage{hyperref}
\usepackage{enumitem}
\usepackage{xcolor}
\usepackage{mdframed}
\usepackage{booktabs}
\usepackage{array}

\geometry{margin=1in}

\definecolor{keycolor}{RGB}{180, 70, 0}
\definecolor{notecolor}{RGB}{240, 240, 255}
\definecolor{takeawaycolor}{RGB}{255, 245, 230}

\newmdenv[
  backgroundcolor=notecolor,
  linecolor=keycolor,
  linewidth=1.5pt,
  topline=false, bottomline=false, rightline=false,
  innerleftmargin=10pt, innerrightmargin=10pt,
]{notebox}

\newmdenv[
  backgroundcolor=takeawaycolor,
  linecolor=keycolor,
  linewidth=2pt,
  roundcorner=5pt,
  innerleftmargin=10pt, innerrightmargin=10pt,
  innertopmargin=8pt, innerbottommargin=8pt,
]{takeawaybox}

\title{
  \textbf{CS 3600: Introduction to Artificial Intelligence}\\[0.5em]
  \large Lecture 9 --- Classification and Machine Learning\\[0.3em]
  \normalsize Aaron Hillegass, Georgia Tech
}
\author{Lecture Notes}
\date{}

\begin{document}

\maketitle
\tableofcontents
\newpage

% ---------------------------------------------------------------
\section{From Regression to Classification}
% ---------------------------------------------------------------

The same ML workflow applies to classification: build a model, define a loss, minimize the loss with gradient descent, then evaluate. The challenge is adapting each piece for the case where the output is a discrete class label rather than a continuous number.

\textbf{Running example:} Female blue whales tend to be longer and heavier than males. Can we classify whale gender from length and weight?

A natural first attempt: use a linear model and predict Female if $\beta_0 + \beta_1 \times \text{length} + \beta_2 \times \text{weight} > 0$. This defines a linear decision boundary. The key difficulty is finding a \emph{differentiable loss function} to optimize the parameters $\beta$.

% ---------------------------------------------------------------
\section{Evaluating a Binary Classifier}
% ---------------------------------------------------------------

\subsection{The Confusion Matrix}

For binary classification, every prediction falls into one of four cells:

\begin{center}
\begin{tabular}{l|cc}
 & \textbf{Predicted Negative} & \textbf{Predicted Positive} \\\hline
\textbf{Actually Negative} & True Negative (TN) & False Positive (FP) \\
\textbf{Actually Positive} & False Negative (FN) & True Positive (TP) \\
\end{tabular}
\end{center}

\subsection{Metrics}

From the confusion matrix we compute four key metrics. Using the stroke example ($n = 345$, TP = 100, TN = 200, FP = 25, FN = 20):

\[
\text{Accuracy} = \frac{\text{TP} + \text{TN}}{n} = \frac{300}{345} \approx 87\%
\]
\[
\text{Recall} = \frac{\text{TP}}{\text{Actual Positives}} = \frac{100}{120} \approx 83\%
\]
\[
\text{Precision} = \frac{\text{TP}}{\text{Predicted Positives}} = \frac{100}{125} \approx 80\%
\]
\[
\text{F1} = 2 \times \frac{\text{Recall} \times \text{Precision}}{\text{Recall} + \text{Precision}} \approx 0.816
\]

\begin{notebox}
\textbf{Recall vs.\ Precision tradeoff:} In medical diagnosis, missing a true positive (low recall) can be fatal --- lower the threshold to catch more positives. In spam filtering, falsely flagging legitimate emails (low precision) is costly --- raise the threshold. The threshold is a business decision, not a model decision.
\end{notebox}

\subsection{ROC Curve and AUC}

Rather than evaluating at a single threshold, a \textbf{ROC} (Receiver Operating Characteristics) curve sweeps the threshold from 0 to 1 and plots the True Positive Rate against the False Positive Rate at each point.

The \textbf{AUC} (Area Under the Curve) summarizes this curve in a single number:
\begin{itemize}
  \item AUC = 1.0: perfect classifier
  \item AUC = 0.5: random classifier (the diagonal)
  \item AUC $<$ 0.5: worse than random (but can be flipped)
\end{itemize}

AUC is threshold-independent and is often the best single metric for comparing classifiers.

% ---------------------------------------------------------------
\section{Loss Functions for Classification}
% ---------------------------------------------------------------

\subsection{Why Not Accuracy?}

Accuracy, recall, and F1 are all \textbf{not differentiable} --- they depend on a hard threshold (predict 1 if $\hat{y} > 0.5$), which creates discontinuities. Gradient descent cannot work with them directly.

\textbf{Solution:} Make the model output a smooth probability-like value, then define a differentiable loss on those probabilities.

\subsection{The Sigmoid (Logistic) Function}

The sigmoid function squashes any real number into $(0, 1)$:
\[
\sigma(x) = \frac{1}{1+e^{-x}}
\]

Useful properties:
\[
1 - \sigma(x) = \sigma(-x) \qquad \sigma^{-1}(p) = \log\!\left(\frac{p}{1-p}\right) \quad \text{for } 0 < p < 1
\]
\[
\sigma'(x) = \frac{e^x}{(1+e^x)^2} = \sigma(x)(1-\sigma(x))
\]

By composing the sigmoid with our linear model, we get soft predictions:
\[
\hat{y} = \sigma(\beta_0 + \beta_1 x_1 + \beta_2 x_2) = \frac{1}{1 + e^{-(\beta_0 + \beta_1 x_1 + \beta_2 x_2)}} \in (0,1)
\]

To make a hard classification, threshold $\hat{y}$ at 0.5 (default), or adjust the threshold to tune recall/precision.

\subsection{Why Not Mean Squared Error for Classification?}

\[
\text{MSE} = \frac{1}{n}\sum_i\left(y_i - \frac{1}{1+e^{-\beta^T x_i}}\right)^2
\]

Two problems:
\begin{enumerate}
  \item MSE with sigmoid is \textbf{not convex}, so gradient descent may converge to a local minimum.
  \item It is conceptually odd: $y$ is always exactly 0 or 1, while $\hat{y} \in (0,1)$.
\end{enumerate}

\subsection{Cross-Entropy Loss}

Derived from maximum likelihood estimation, cross-entropy is the correct loss for classification. For binary labels $y_i \in \{0,1\}$:
\[
\boxed{\text{CE} = -\frac{1}{n}\sum_i\bigl[y_i\log(\hat{y}_i) + (1-y_i)\log(1-\hat{y}_i)\bigr]}
\]

Why this works:
\begin{itemize}
  \item When $y_i = 1$ and $\hat{y}_i \to 1$: loss $\to 0$ (correct, confident).
  \item When $y_i = 1$ and $\hat{y}_i \to 0$: loss $\to \infty$ (wrong and confident --- heavily penalized).
  \item Makes the optimization problem \textbf{convex}, guaranteeing a global minimum.
\end{itemize}

\begin{takeawaybox}
\textbf{Logistic Regression} is simply: (1) linear combination of features $\beta^T x$, (2) feed through sigmoid, (3) optimize $\beta$ with cross-entropy loss. Despite the name, it is a \emph{classifier}, not a regressor.
\end{takeawaybox}

\subsection{Regularization: Preventing Runaway Parameters}

If the training data is linearly separable, cross-entropy loss pushes the parameters toward $\pm\infty$, turning the sigmoid into a step function. This model is overconfident and generalizes poorly.

To prevent this, add an L2 regularization penalty:
\[
\mathcal{L}(\beta) = -\frac{1}{n}\sum_i\bigl[y_i\log(\hat{y}_i) + (1-y_i)\log(1-\hat{y}_i)\bigr] + \lambda(\beta^T\beta)
\]
The hyperparameter $\lambda > 0$ controls the strength of regularization, penalizing large weights and keeping the model calibrated.

% ---------------------------------------------------------------
\section{Multi-Class Classification}
% ---------------------------------------------------------------

\subsection{Extending to $k > 2$ Classes}

For $k$ classes, we assign each class its own parameter vector $\vec{\beta}_c$. For each input $\vec{x}$, we compute a score $z_c = \vec{\beta}_c^T\vec{x}$ for every class. A hard classifier would simply pick $\arg\max_c z_c$, but \texttt{max} is not differentiable.

\subsection{Softmax: A Differentiable Argmax}

The \textbf{LogSumExp} function, $\log\sum_k e^{x_k}$, is a smooth approximation to \texttt{max}: when one $x_k$ dominates, $\log\sum e^{x_i} \approx x_k$. Its gradient is the \textbf{Softmax}:
\[
\text{softmax}(z_i) = \frac{e^{z_i}}{\sum_j e^{z_j}}
\]

Properties:
\begin{itemize}
  \item $0 < \text{softmax}(z_i) < 1$ for all $i$
  \item $\sum_j \text{softmax}(z_j) = 1$ (outputs are a probability distribution)
  \item Differentiable everywhere
\end{itemize}

\textbf{Numerical example:} Given scores $\boldsymbol{Z} = [-1.9,\ 2.1,\ 4.8]$:
\[
e^{\boldsymbol{Z}} = [0.1496,\ 8.166,\ 121.51], \quad \Sigma(e^{\boldsymbol{Z}}) = 129.82
\]
\[
\text{softmax}(\boldsymbol{Z}) = [0.00115,\ 0.0629,\ 0.9359]
\]
The model is 93.6\% confident in class 3.

\subsection{Temperature-Scaled Softmax}

The \textbf{temperature} parameter $T$ controls how ``peaked'' the distribution is:
\[
\text{softmax}_T(r_i) = \frac{\exp(r_i/T)}{\sum_j \exp(r_j/T)}
\]
\begin{itemize}
  \item $T \to 0$: approaches a one-hot (argmax) distribution.
  \item $T = 1$: standard softmax.
  \item $T \to \infty$: approaches a uniform distribution (maximum uncertainty).
\end{itemize}
Temperature scaling is used in knowledge distillation and controlling the randomness of language model outputs.

\subsection{Evaluating Multi-Class Classifiers}

For $k$ classes, the confusion matrix generalizes to a $k \times k$ table. Per-class precision, recall, and F1 can be computed by treating each class as a binary ``one vs.\ rest'' problem, then averaged (macro or weighted) for a single summary score.

% ---------------------------------------------------------------
\section{Multi-Class Loss: Generalizing Cross-Entropy}
% ---------------------------------------------------------------

\subsection{One-Hot Encoding and Matrix Formulation}

For $n$ samples and $k$ classes:
\begin{itemize}
  \item $\boldsymbol{Y}$ is an $n \times k$ \textbf{one-hot} matrix: $y_{i,j} = 1$ if sample $i$ is of class $j$, else 0.
  \item $\hat{\boldsymbol{Y}}$ is an $n \times k$ matrix of softmax outputs: $\hat{y}_{i,j}$ is the predicted probability that sample $i$ is of class $j$.
\end{itemize}

The \textbf{Mean Cross-Entropy} loss generalizes directly:
\[
\boxed{\mathcal{L}(\beta) = -\sum_{i=1}^{n}\left[\sum_{j=1}^{k} y_{i,j}\log(\hat{y}_{i,j})\right]}
\]

Because $y_{i,j}$ is one-hot, only the log-probability of the correct class contributes to the loss for each sample.

\subsection{Matrix Form of the Full Model}

With $d$ input features, $n$ samples, and $k$ classes:
\[
\boldsymbol{X}_{n\times(d+1)},\quad \boldsymbol{\beta}_{(d+1)\times k},\quad \boldsymbol{Z} = \boldsymbol{X}\boldsymbol{\beta},\quad \hat{\boldsymbol{Y}} = \text{softmax}(\boldsymbol{Z}),\quad \mathcal{L} = \text{cross-entropy}(\hat{\boldsymbol{Y}},\boldsymbol{Y})
\]

\subsection{Gradient via the Chain Rule}

A beautiful result: the combined gradient of softmax + cross-entropy with respect to the pre-softmax scores $\boldsymbol{Z}$ is simply:
\[
\frac{\partial\mathcal{L}}{\partial z_{i,j}} = \hat{y}_{i,j} - y_{i,j}
\]
This is prediction minus truth --- elegant and computationally cheap. Applying the chain rule through $\boldsymbol{Z} = \boldsymbol{X}\boldsymbol{\beta}$:
\[
\boxed{\nabla_\beta \mathcal{L} = \boldsymbol{X}^T(\hat{\boldsymbol{Y}} - \boldsymbol{Y})}
\]

This compact formula drives parameter updates for the entire multi-class logistic regression model. Note that one class's parameter vector is redundant (can be fixed to zero) since softmax scores are relative --- this is \textbf{multinomial logistic regression}.

% ---------------------------------------------------------------
\section{Summary and Key Takeaways}
% ---------------------------------------------------------------

\begin{takeawaybox}
\textbf{Key Takeaways from Lecture 9}
\begin{enumerate}
  \item \textbf{Accuracy/F1/Recall are not differentiable} --- we cannot use them directly as loss functions. We need smooth probability outputs.

  \item \textbf{Sigmoid} squashes any real value into $(0,1)$. Combined with a linear model, it gives logistic regression.

  \item \textbf{Cross-entropy loss} is the correct loss for classification. It is convex (for logistic regression), and heavily penalizes confident wrong predictions.

  \item \textbf{L2 regularization} $(\lambda\beta^T\beta)$ prevents parameters from exploding when data is linearly separable.

  \item \textbf{Softmax} generalizes sigmoid to $k > 2$ classes. It is the differentiable approximation to argmax, with outputs summing to 1.

  \item The combined gradient of \textbf{softmax + cross-entropy} is just $\hat{Y} - Y$ --- prediction minus truth. The full gradient is $\nabla_\beta\mathcal{L} = X^T(\hat{Y} - Y)$.

  \item The \textbf{ROC curve / AUC} evaluates a classifier across all thresholds, giving a threshold-independent quality measure.

  \item \textbf{Temperature} controls softmax sharpness: low $T$ = confident, high $T$ = uniform.
\end{enumerate}
\end{takeawaybox}

\subsection{Further Resources}
\begin{itemize}
  \item \textbf{Slides:} by Aaron Hillegass, Georgia Tech
  \item \textbf{Proof of softmax/CE gradient:} \url{https://medium.com/analytics-vidhya/derivative-of-log-loss-function-for-logistic-regression-9b832f025c2d}
\end{itemize}

\end{document}
